\begin{abstract}
Autonomous vehicles increasingly run computation-heavy tasks with strict service-quality demands that their onboard processors cannot always meet on their own, which is why offloading part of that work to the network edge has become common practice. When many vehicles compete for the same wireless link and the same edge server, however, a policy that offloads aggressively without regard for channel conditions or queue congestion can fail badly, and this failure mode is easy to miss if a system is only ever tested at one vehicle density. This paper presents MA-ROPE, a Multi-RAT (Radio Access Technology) Privacy-Preserving Federated Deep Reinforcement Learning framework for Vehicle-to-Everything (V2X) task offloading. MA-ROPE jointly chooses among four radio interfaces (LTE, ITS-G5, Wi-Fi 6, and 5G-NR) and decides how much of each task to offload, using a Dueling Double Deep Q-Network (D3QN) whose training is coordinated across vehicles through Federated Averaging (FedAvg) so that raw training observations are not pooled across vehicles (the offloaded share of a task is still sent to the RSU). Under dense, latency-critical traffic, MA-ROPE significantly improves the rate at which tasks are completed on time compared to fixed single-radio baselines. We also introduce the Deadline Retention Ratio, a new metric for how much of a policy's on-time performance survives a large increase in fleet size, and find that MA-ROPE holds onto nearly all of its performance while baseline methods experience severe degradation. Furthermore, the joint radio-selection and offloading logic adds no measurable computational overhead over simpler controllers. Ultimately, radio selection, partial offloading, and edge-queue contention must be optimized jointly, not separately, for vehicular networks to deliver time-critical tasks reliably as they scale.
\end{abstract}
\begin{IEEEkeywords}
Task offloading, Vehicle-to-Everything (V2X) communication, vehicular edge computing, federated deep reinforcement learning, multi-RAT selection, heterogeneous networks, edge computing, autonomous driving, privacy-preserving learning, Jakes fading.
\end{IEEEkeywords}
